\section{Limitations and Future Work}

Several limitations bound the strength of the present claims, and each points to a concrete next step.

The most serious is that the learnable-drug subset was selected using all \NLines{} cell lines, including
the validation and test lines, so any number reported on that subset is a best-case diagnostic rather than
a generalization estimate; the selection should instead be performed inside each cross-validation fold,
using only the training lines. The subset is moreover chosen from label statistics alone and was not
validated in this report against the per-drug performance the model actually achieves, so it should be
read as a diagnostic device rather than a claim that only ten drugs are learnable.

The scGPT advantage over PCA is small: a Spearman gap of $+0.036$ at K\,$=$\,\NLearn{} and $+0.012$ at
K\,$=$\,\NDrugs{}, of the same order as the seed-to-seed variation. Across three seeds the gap averages
$+0.04$ and is not sign-consistent, ranging from $-0.00$ to $+0.09$, and it does not survive DrEval's
normalized metric, so it should not be presented as an established margin until it has been reproduced over
more seeds and a wider drug set. The external benchmark is itself preliminary, covering only the
leave-cell-line-out split on the learnable subset rather than the leave-tissue-out and leave-drug-out
settings or the full panel.

The clearest opportunities for improvement follow from the diagnosis that the task is label-limited. First,
the training target can be aligned with the evaluation by regressing on the doubly normalized residual, in
which both the drug and the cell-line mean effects are removed, with the means estimated only on the
training lines of each fold; this matches the quantity DrEval rewards, and an initial check
(\code{outputs/dreval/dreval\_normalized.csv}) indicates that about a fifth of the current signal is the
cell-line effect alone. Second, the single-cell resolution has
yet to justify itself, since averaging a line's cells loses nothing measurable; a pooling model that
predicts the line label from a bag of cells, for example through attention, is the natural test and must
beat the ridge control to be worthwhile. Third, and most directly, the ceiling can be raised by adding
independent cell lines, as CTRPv2 alone screens roughly 1{,}100 against the \NLines{} in the current
overlap. Beyond these, a diagnostic error analysis of where the model fails is low-risk and immediately
useful, whereas gene-level interpretation to surface transcriptomic drivers of resistance is worthwhile
only once per-drug performance is substantially higher and stable, so that the interpretation does not
merely describe noise. The longer horizon remains a cross-database, multi-task model spanning
efficacy and toxicity, fine-tunable on clinical outcomes.
